\documentclass[aspectratio=169,12pt]{beamer}

\usetheme{metropolis}
\usepackage[utf8]{inputenc}
\usepackage[spanish]{babel}
\usepackage{graphicx}
\usepackage{booktabs}
\usepackage{amssymb}
\usepackage{xcolor}
\usepackage{tikz}
\usetikzlibrary{positioning, arrows.meta}

% Figuras compartidas con el reporte
\graphicspath{{../figuras/}}
\newcommand{\figdir}{../figuras}
% Imagen con placeholder automático (compila con o sin el archivo)
\newcommand{\imgph}[2][0.85]{%
  \IfFileExists{\figdir/#2}{\includegraphics[width=#1\linewidth]{#2}}%
  {\fcolorbox{gray}{gray!12}{\begin{minipage}[c][3.0cm][c]{#1\linewidth}%
     \centering\ttfamily\footnotesize [PLACEHOLDER]\\[3pt]\detokenize{#2}%
   \end{minipage}}}%
}

\title{El Sistema de Adquisición de Datos (DAQ) de Alta Velocidad}
\subtitle{Integración y validación con ADC físico AD9226}
\author{Eduardo Santos Mena}
\institute{Asesores: Dr. Samuel Pérez Huerta · Dr. Augusto David Ariza Flores\\
Unidad Académica de Ciencia y Tecnología de la Luz y la Materia, UAZ}
\date{Avance — Cuarto Semestre 2026} % TODO: confirmar fecha

\begin{document}

\maketitle

% ---------------------------------------------------------------------------
\begin{frame}{Contexto}
  \textbf{Tomografía fotoacústica (PAT) + IA} para la \textbf{detección temprana de
  artritis reumatoide} (AR) — proyecto de doctorado (LUMAT–UAZ).
  \begin{itemize}
    \item Excitación óptica + detección ultrasónica, sin radiación ionizante.
    \item La PAT revela cambios vasculares precoces (angiogénesis sinovial).
    \item Metas del sistema: resolución 150–200 $\mu$m, profundidad 1–2 cm.
  \end{itemize}
  \begin{block}{Estrategia}
    (a) instrumentación de adquisición propia de bajo costo \quad
    (b) algoritmos (CNN) de reconstrucción/segmentación.
  \end{block}
\end{frame}

% ---------------------------------------------------------------------------
\begin{frame}{Continuidad del proyecto}
  \begin{itemize}
    \item \textbf{Sem 2:} fuente de excitación — driver de avalancha, láser LED
      445 nm, pulsos $<$30 ns, bajo costo.
    \item \textbf{Sem 3:} Driver V2.0 + núcleo digital validado con \textit{Virtual
      ADC} (señales sintéticas).
    \item \textbf{Sem 4 (este):} \alert{integración del ADC físico y validación del
      DAQ de alta velocidad}; primera caracterización de la detección.
  \end{itemize}
\end{frame}

% ---------------------------------------------------------------------------
\begin{frame}{Objetivo del semestre}
  \begin{alertblock}{Objetivo central}
    Validar la cadena de adquisición de datos (DAQ) con \textbf{hardware físico}, a
    la velocidad y resolución de diseño.
  \end{alertblock}
  \begin{itemize}
    \item Integración física ADC (AD9226) — FPGA.
    \item Operación del Driver V2.0 y estimación de potencia óptica.
    \item Caracterización inicial de la detección (\textit{etapa crítica anticipada}).
  \end{itemize}
\end{frame}

% ---------------------------------------------------------------------------
\begin{frame}{Arquitectura del DAQ — etapas y subetapas}
\centering
\resizebox{\linewidth}{!}{%
\begin{tikzpicture}[
  font=\sffamily,
  block/.style={rectangle, draw=blue!55, thick, rounded corners, fill=blue!5,
                align=center, text width=3.1cm, minimum height=2.4cm, inner sep=4pt},
  arr/.style={-{Stealth[length=3mm]}, very thick, draw=blue!55}
]
\node[block] (exc) {\textbf{1. Excitación}\\[3pt]\scriptsize Driver V2.0\\
  \scriptsize trigger 5\,kHz\\ \scriptsize monitor óptico (TIA)};
\node[block, right=0.7cm of exc] (trg) {\textbf{2. Trigger}\\[3pt]\scriptsize pin 77 / S2\\
  \scriptsize detección de flanco\\ \scriptsize arranque de ventana};
\node[block, right=0.7cm of trg] (adc) {\textbf{3. Conversión}\\[3pt]\scriptsize AD9226 $\cdot$ 12 bit\\
  \scriptsize 54 MSPS (PLL $\times$2)\\ \scriptsize $\pm$5 V, LSB 2.44 mV};
\node[block, right=0.7cm of adc] (cap) {\textbf{4. Captura}\\[3pt]\scriptsize store-and-forward\\
  \scriptsize BRAM 270 muestras\\ \scriptsize ventana 5\,$\mu$s};
\node[block, right=0.7cm of cap] (uart) {\textbf{5. Enlace}\\[3pt]\scriptsize UART\\
  \scriptsize header sync 4 B\\ \scriptsize FW version};
\node[block, right=0.7cm of uart] (host) {\textbf{6. Host}\\[3pt]\scriptsize Python\\
  \scriptsize window\_live\\ \scriptsize promediado};
\draw[arr] (exc)--(trg); \draw[arr] (trg)--(adc); \draw[arr] (adc)--(cap);
\draw[arr] (cap)--(uart); \draw[arr] (uart)--(host);
\end{tikzpicture}%
}
\vfill
{\footnotesize Estrategia \textit{store-and-forward}: captura rápida en BRAM,
transmisión en el tiempo muerto entre disparos.}
\end{frame}

% ---------------------------------------------------------------------------
\begin{frame}{Validación incremental (\textit{bring-up})}
\centering\scriptsize
\begin{tabular}{cll}
\toprule
\textbf{Etapa} & \textbf{Subetapa que valida} & \textbf{Estado} \\
\midrule
0 & Reloj + configuración del bitstream & \checkmark \\
1 & Trigger (S2 / pin 77) y flanco & \checkmark \\
2 & Enlace UART aislado & \checkmark \\
3 & Lectura del ADC (voltímetro) & \checkmark \\
4 & Captura por reloj interno (5 kHz) & \checkmark \\
5 & Captura por trigger externo & \checkmark \\
6 & Ráfaga 270 muestras @ 27 MSPS & \checkmark \\
7 & Ráfaga @ 54 MSPS (PLL $\times$2) & \checkmark \\
8 & Ráfaga @ 54 MSPS + 12 bits & \checkmark \\
\bottomrule
\end{tabular}
\vfill
{\footnotesize Cada etapa prueba una sola variable; no se avanza hasta que pasa.}
\end{frame}

% ---------------------------------------------------------------------------
\begin{frame}{Algoritmo de captura: máquina de estados}
\centering
\resizebox{0.95\linewidth}{!}{%
\begin{tikzpicture}[
  font=\sffamily,
  st/.style={rectangle, draw=blue!55, thick, rounded corners, fill=blue!5,
             minimum width=1.8cm, minimum height=0.95cm, align=center},
  tr/.style={-{Stealth[length=2.5mm]}, thick, draw=blue!55}
]
\node[st] (idle) {IDLE};
\node[st, right=1.5cm of idle] (cap) {CAPTURE};
\node[st, right=1.5cm of cap] (prep) {PREP};
\node[st, right=1.5cm of prep] (hdr) {HDR};
\node[st, right=1.5cm of hdr] (send) {SEND};
\draw[tr] (idle) -- node[above,font=\scriptsize]{trig} (cap);
\draw[tr] (cap) -- node[above,font=\scriptsize]{ptr=N-1} (prep);
\draw[tr] (prep) -- (hdr);
\draw[tr] (hdr) -- node[above,font=\scriptsize]{hdr$\times$4} (send);
\draw[tr] (send) to[out=-90,in=-90] node[below,font=\scriptsize]{ptr=N} (idle);
\end{tikzpicture}%
}
\vfill
{\footnotesize \textit{Store-and-forward}: IDLE vigila el trigger $\rightarrow$
CAPTURE llena la BRAM $\rightarrow$ PREP/HDR/SEND transmiten. Fuera de IDLE, los
triggers se ignoran $\Rightarrow$ ráfagas íntegras.}
\end{frame}

% ---------------------------------------------------------------------------
\begin{frame}{Sincronización: captura rápida, transmisión lenta}
\centering
\imgph[0.9]{fig-timing-fsm.png}\\[3pt]
{\footnotesize La transmisión ($\sim$ms) excede el periodo del trigger (200 µs)
$\Rightarrow$ se omiten disparos $\Rightarrow$ \alert{promediado coherente}
(SNR $\propto \sqrt{N}$).}
\end{frame}

% ---------------------------------------------------------------------------
\begin{frame}{Tasa de muestreo: PLL ×2 → 54 MSPS}
\begin{columns}[T]
\begin{column}{0.56\textwidth}
  \imgph[1.0]{fig-muestreo-nyquist.png}
\end{column}
\begin{column}{0.44\textwidth}
  \begin{itemize}
    \item 27 MHz $\rightarrow$ 54 MSPS (Ts: 37 $\rightarrow$ 18.5 ns).
    \item 2.5 MHz: 10.8 $\rightarrow$ \alert{21.6 muestras/ciclo}.
    \item Sobremuestreo holgado del \textit{ringing}.
  \end{itemize}
\end{column}
\end{columns}
\end{frame}

% ---------------------------------------------------------------------------
\begin{frame}{Profundidad de codificación: 8 → 12 bits}
\begin{columns}[T]
\begin{column}{0.56\textwidth}
  \imgph[1.0]{fig-niveles-codificacion.png}
\end{column}
\begin{column}{0.44\textwidth}
  \begin{itemize}
    \item 256 $\rightarrow$ \alert{4096 niveles}.
    \item LSB 39 mV $\rightarrow$ \alert{2.44 mV} ($\times$16).
    \item Clave para señales acústicas débiles.
  \end{itemize}
\end{column}
\end{columns}
\end{frame}

% ---------------------------------------------------------------------------
\begin{frame}{Memoria: el BSRAM no es el límite}
\begin{columns}[T]
\begin{column}{0.58\textwidth}
  \imgph[1.0]{fig-memoria.png}
\end{column}
\begin{column}{0.42\textwidth}
  \begin{itemize}
    \item Memoria $= N \times b$ bits.
    \item Ventana útil ($\sim$24 Kbit) $\ll$ BSRAM ($\sim$468 Kbit).
    \item El límite real es la \alert{UART}, no la memoria.
  \end{itemize}
\end{column}
\end{columns}
\end{frame}

% ---------------------------------------------------------------------------
\begin{frame}{Resultado central: DAQ validado}
\begin{columns}[T]
\begin{column}{0.5\textwidth}
  \begin{itemize}
    \item Captura, almacena y transmite \alert{sin pérdida de datos}.
    \item \textbf{54 MSPS}, \textbf{12 bits} (LSB 2.44 mV).
    \item $\sim$21.6 muestras/ciclo a 2.5 MHz $\rightarrow$ resuelve el
      \textit{ringing} del sensor.
    \item Supera lo proyectado para el periodo.
  \end{itemize}
\end{column}
\begin{column}{0.5\textwidth}
  \centering
  \imgph[1.0]{daq-captura.png}\\[2pt]
  {\scriptsize Captura a la velocidad y niveles especificados.}
\end{column}
\end{columns}
\end{frame}

% ---------------------------------------------------------------------------
\begin{frame}{Driver V2.0 y estimación de potencia}
\begin{columns}[T]
\begin{column}{0.5\textwidth}
  \centering \imgph[1.0]{pulsos-corriente.png}\\[2pt]
  {\scriptsize Corriente del diodo (shunt 0.1 $\Omega$).}
\end{column}
\begin{column}{0.5\textwidth}
  \centering \imgph[1.0]{potencia-laser.png}\\[2pt]
  {\scriptsize Estimación de potencia V$\cdot$I (preliminar).}
\end{column}
\end{columns}
\end{frame}

% ---------------------------------------------------------------------------
\begin{frame}{La etapa crítica anticipada: la detección}
  Convertir una onda acústica \textbf{débil} ($\mu$V–mV), captada por un transductor
  de \textbf{alta impedancia}, en señal observable —y en presencia de una excitación
  de \textbf{alto voltaje}— se anticipó como la parte \alert{más difícil} del sistema.
  \begin{block}{Primer intento}
    Con el DAQ ya validado, se llevó el sistema al banco para la primera adquisición.
  \end{block}
\end{frame}

% ---------------------------------------------------------------------------
\begin{frame}{Caracterización: la señal observada es \textit{pickup}}
\begin{columns}[T]
\begin{column}{0.54\textwidth}
  CH2 mostró $\sim$150 mV pp coincidentes con el láser ($\sim$2.6 MHz).
  \begin{alertblock}{Prueba decisiva}
    \textbf{Desconexión del sensor:} la onda \emph{persiste} con el transductor
    desconectado $\Rightarrow$ no proviene del piezo. Por sí sola basta.
  \end{alertblock}
  Dos controles que \textbf{corroboran}:
  \begin{itemize}
    \item \textbf{$t=0$:} sin retardo de tiempo de vuelo $\rightarrow$ es
      electromagnética, no acústica.
    \item \textbf{Distancia:} no cambia al alejar la fuente (sin dependencia de TOF).
  \end{itemize}
\end{column}
\begin{column}{0.46\textwidth}
  \centering \imgph[1.0]{ringing-falso.png}\\[2pt]
  {\scriptsize Control de descarte (osciloscopio).}
\end{column}
\end{columns}
\end{frame}

% ---------------------------------------------------------------------------
\begin{frame}{Evidencia en video: ruido EM del disparo}
  \centering
  \vfill
  {\Large\textbf{[ VIDEO ]}}\\[8pt]
  Ruido electromagnético sincronizado con el disparo de alto voltaje.\\[6pt]
  {\footnotesize Reproducir el clip durante la presentación.}\\[2pt]
  {\scriptsize \texttt{evidencia/fotos/evidencia\_ruido\_trigger.mp4}}
  \vfill
\end{frame}

% ---------------------------------------------------------------------------
\begin{frame}{Por qué: el front-end de detección}
  \begin{itemize}
    \item El piezo (alta-Z) va \textbf{en crudo} al digitalizador, \textbf{sin
      preamplificador}.
    \item La señal PA ($\mu$V–mV) queda bajo el piso de observación y bajo el
      \textit{pickup} de 150 mV.
    \item El subsistema digital (ADC, FPGA, enlace, Driver) está validado: el reto
      está \alert{aislado} en la interfaz analógica.
  \end{itemize}
  \begin{exampleblock}{Valor del resultado}
    No es un obstáculo inesperado: \textbf{localiza} con precisión el trabajo
    restante y reduce el riesgo de la siguiente fase.
  \end{exampleblock}
\end{frame}

% ---------------------------------------------------------------------------
\begin{frame}{Acoplamiento de impedancias → preamplificador}
\begin{columns}[T]
\begin{column}{0.56\textwidth}
  \imgph[1.0]{fig-impedancia.png}
\end{column}
\begin{column}{0.44\textwidth}
  \begin{itemize}
    \item Piezo = fuente de \textbf{alta-Z}.
    \item Entrada de 50 $\Omega$ \alert{atenúa} la señal; 1 M$\Omega$ la preserva.
    \item Se necesita \textbf{preamp de alta-Z y bajo ruido, junto al sensor}.
    \item Su salida de baja-Z $\rightarrow$ inmune al \textit{pickup} en el cable.
  \end{itemize}
\end{column}
\end{columns}
\end{frame}

% ---------------------------------------------------------------------------
\begin{frame}{Trabajo futuro: cadena de detección}
  \begin{itemize}
    \item \textbf{Bring-up de detección} (D0–D4): salud del piezo $\rightarrow$
      toque mecánico $\rightarrow$ fuente acústica acoplada $\rightarrow$ preamp
      $\rightarrow$ láser.
    \item \textbf{Preamplificador} de bajo ruido (AD8331/AD8332): alta-Z, 40–60 dB,
      0.1–5 MHz, $\leq$3–5 nV/$\sqrt{\text{Hz}}$.
    \item \textbf{Control del EMI}: disparo óptico (t=0 limpio), apantallamiento,
      ventana retardada.
    \item \textbf{Ampliar ventana} (FIFO) $\rightarrow$ A-scans en fantomas
      $\rightarrow$ reconstrucción + IA.
  \end{itemize}
\end{frame}

% ---------------------------------------------------------------------------
\begin{frame}{Cronograma del siguiente periodo}
\centering\scriptsize
\begin{tabular}{cll}
\toprule
\textbf{Bloque} & \textbf{Actividad} & \textbf{Entregable} \\
\midrule
B1  & Diseño del front-end (preamp, Z, shielding, sincronía) & Esquemático + ruido \\
B1' & Validación de recepción (Plan B-1/B-2) & Señal acústica real \\
B2  & Fabricación y prueba del front-end & Cadena caracterizada \\
B3  & Firmware: ventana ampliada (FIFO) & A-scan útil \\
B4  & Adquisición PA en fantoma & A-scan con TOF coherente \\
B5  & Procesamiento (retroproyección) & Reconstrucción básica \\
B6  & IA (denoising) + escritura & Artículo / reporte \\
\bottomrule
\end{tabular}
\end{frame}

% ---------------------------------------------------------------------------
\begin{frame}{Conclusiones}
  \begin{enumerate}
    \item \textbf{DAQ validado (avance central):} 12 bits, 54 MSPS, sin pérdida de
      datos, con hardware físico.
    \item \textbf{Driver V2.0 y potencia:} operación estable + estimación por pulso.
    \item \textbf{Detección localizada:} la lectura es \textit{pickup}, no acústica;
      el front-end analógico define el trabajo futuro.
  \end{enumerate}
  \begin{center}
    \alert{Plataforma de adquisición robusta + camino claro hacia la imagen PA.}
  \end{center}
\end{frame}

% ---------------------------------------------------------------------------
\begin{frame}[standout]
  Gracias.\\[4pt]
  \small Preguntas
\end{frame}

\end{document}
